\subsection{Abelian Selmer fibres over rational compact points}

\begin{statement}{A.64}{The exact abelian Selmer fibre over a rational compact point}{Proved}
Let $E/k$ be any elliptic curve over a number field and let
$D\subset E$ be a reduced divisor with at least two geometric
points. Put $U=E-D$. In the full geometric abelian quotient of
$\pi_1(U)$, the Selmer fibre over the compact section of a point
$B\in E(k)$ is as follows:
\[
 \begin{array}{c|c}
 B\in U(k)&\text{the single point section of }B\\
 B\in D(k)&\text{the complete cusp packet at }B,
          \text{ parametrized injectively by }H^1(k,\widehat\Z(1)).
 \end{array}
 \tag{A64.1}
\]
The packet parametrization uses a chosen nonzero $k$-rational
tangent vector. Only local geometricity at the finite places is
needed for this assertion. In particular A.63 gives the complete
abelian Selmer fibres on all finite alterations of integer
exponents: the point section at $nR$ for $n\ne0$, and the whole
$O$-cusp packet for $n=0$.
\end{statement}
\begin{proof}
We first prove the cohomological fact, including the prime two.
For every number field $F$, localization at all finite places is
injective on $H^1(F,\widehat\Z(1))$, as also established in B.26.
Here is a direct proof. For odd $\ell$ and finite level $\ell^r$,
restrict a locally zero class to $F(\mu_{\ell^r})$. The coefficient
module is then constant, so a nonzero class would be a nonzero
finite character, detected by Chebotarev at a finite place.
The original class consequently inflates from
$H^1(\operatorname{Gal}(F(\mu_{\ell^r})/F),\mu_{\ell^r})$.
That Galois group is cyclic. Choose an unramified finite place
whose Frobenius generates it. Inflation into the local $H^1$
is injective, so local triviality kills this class also. This
proves injectivity at every odd prime level.

For $\ell=2$, first pass to $F(i)$. Every group
$\operatorname{Gal}(F(i,\mu_{2^r})/F(i))$ is cyclic, being a
subgroup of $1+4\Z/2^r\Z$, so the same argument applies there.
Passing to the inverse limit, a locally zero class in
$H^1(F,\Z_2(1))$ restricts to zero over $F(i)$.
Restriction--corestriction makes it $2$-torsion. The exact
sequence
\[
 0\longrightarrow\Z_2(1)\longrightarrow\Q_2(1)
   \longrightarrow\Q_2/\Z_2(1)\longrightarrow0
\]
identifies the torsion subgroup of this $H^1$ with
$\mu_{2^\infty}(F)$. It injects into the corresponding local
group at a dyadic place, since those roots of unity inject and
$H^0(F_v,\Q_2(1))=0$. Thus the locally zero class is zero.
The finite-level invariant groups are finite, so there is no
$\varprojlim^1$ ambiguity in these $H^1$ limits. Taking all
coefficient primes proves the asserted completed injection.

More generally, let $T$ be a finite-rank compact Galois lattice
which becomes a direct sum of copies of $\widehat\Z(1)$ over
some finite Galois extension $L/k$. If a class in $H^1(k,T)$
is locally zero at all finite places, its restriction to $L$
is zero by the preceding result applied to every component.
Moreover $T^{G_L}=0$, because a number field has only finitely
many roots of unity. Inflation--restriction therefore makes
$H^1(k,T)\to H^1(L,T)$ injective. We have proved
\[
 H^1(k,T)\lhook\joinrel\longrightarrow
       \prod_{v\text{ finite}}H^1(k_v,T).
 \tag{A64.2}
\]

The geometric abelian module $V$ of $U$ sits in the exact
sequence
\[
 0\longrightarrow W_D\longrightarrow V\longrightarrow T E
       \longrightarrow0,\qquad
 W_D=\widehat\Z(1)[D(\overline k)]/
             \widehat\Z(1)\!\left(\sum_{d\in D(\overline k)}d\right).
 \tag{A64.3}
\]
The first term is a lattice of the type in (A64.2).
Choose the actual point section at $B$ if $B\in U$, or a
tangential cusp section at $B$ otherwise. This is a global
reference splitting. A section with the same compact class
differs from it by a class $\xi\in H^1(k,V)$ mapping to zero
in $H^1(k,T E)$, hence by a unique class
$\eta\in H^1(k,W_D)$. Uniqueness uses
$H^0(k,T E)=0$. The same invariant group vanishes at every
finite completion, since the local torsion subgroup of an
elliptic curve is finite; consequently
$H^1(k_v,W_D)\to H^1(k_v,V)$ is injective.

At every finite place the full compact Kummer map is injective
on $E(k_v)$. Thus local Selmerness forces the compact centre
to be $B$. If $B\in U$, the local section is the point section
of $B$, so every localization of $\eta$ vanishes. Apply
(A64.2) with $T=W_D$ to obtain $\eta=0$, proving the first
line of (A64.1).

If $B\in D(k)$, its inertia image
$I_B=\widehat\Z(1)\,B\subset W_D$ is a genuine rank-one
sublattice, because $D$ has at least two geometric points.
The quotient is
\[
 W_D/I_B\simeq
 \widehat\Z(1)[(D-B)(\overline k)]/
       \widehat\Z(1)\!\left(\sum_{d\ne B}d\right),
\]
again a lattice splitting into cyclotomic copies. Locally the
section is a cusp section centred at $B$, so the image of
$\eta$ in $H^1(k,W_D/I_B)$ is locally zero. Equation (A64.2)
kills that image globally. Exactness then places $\eta$ in the
image of $H^1(k,I_B)$, exactly the tangent Kummer packet.
Finally $H^0(k,W_D/I_B)=0$, so this packet maps injectively
to $H^1(k,W_D)$ and hence to section classes. This argument
uses finite completions throughout and asserts no disjoint
local parametrization at the real places.
\end{proof}
